% !TeX root = ../papers/nc-p3.tex
\section{二変数形式の積への商と非可逆境界の排除}
\label{sec:binary-product-obstruction}
$P_j=H^0(\mathbb P^1,\mathcal O(j))$と置く。
\begin{lemma}[二変数形式の積の双対関係]\label{lem:binary-diagonal-paths}
一次生成空間を$P_d^*,P_e^*$と交互に置き、二次関係を乗法
$P_d\otimes P_e\to P_{d+e}$の双対の像とその転置に取った経路代数を考える。
長さ$n\ge2$の次数列を$d_1,\ldots,d_n$、端の位置の$v_i$を1、
内部の$v_i$を2とすると、経路空間の次元は
\[
 \prod_{i=1}^n\max(0,d_i-v_i+1)
\]
である。特に一方の次数が2、他方が3なら指数成長する。
\end{lemma}
\begin{proof}
自由経路空間の双対を、各位置に別の斉次座標$(s_i,t_i)$を用いた多重斉次形式と同一視する。
位置$i,i+1$に関係を挿入した空間の直交補は、この二変数組を同一視したとき零となる形式である。
対角の斉次イデアルは$s_i t_{i+1}-t_i s_{i+1}$で生成される。
全位置の条件を同時に課すと、多項式環の相異なる既約元であるこれらの式の積で割り切れる。
逆に積の任意の倍数は条件を満たす。残る多重次数が$d_i-v_i$なので式が従う。
\end{proof}

\begin{theorem}[全ての導手鉛筆の指数成長]\label{thm:all-conductor-pencils-excluded}
定理\ref{thm:integral-apolar-alternatives}の非可逆模型では、節点・尖点のどちらでも
任意の基点自由鉛筆の反転二窓は指数成長する。従って標準ASではない。
片側の極限二次式が整な四次完全交叉を切る標準AS特殊化は、
その同じ曲線上の二つの可逆層と基点自由鉛筆から得られる。
\end{theorem}
\begin{proof}
導手因子$\Delta$を切る二次形式を$h$とする。$A$に含まれる三次元空間$hP_2$に双対で射影し、
もう一方は$B=P_3$のままとする。六関係はもともと$P_7/W$上の汎関数を積で引き戻した像なので、
射影後は$P_2\otimes P_3\to P_5$の双対像に含まれる。
従って補題の代数への次数付き全射がある。実際、$W$は例外直線の全一次切断に同型に制限され、
$hP_5$はその上で零なので、$hP_5\to P_7/W$は六次元同士の同型であり、射影後の関係は双対像全体である。
$B$始点の長さ$2m+1$で商の次元は$9\cdot 2^{m-1}$ ($m\ge1$)。
$m=8$で既に$1152>\binom{20}{3}=1140$である。滑らかな場合を含む最後の主張は
命題\ref{prop:one-smooth-quartic-specialization}と定理\ref{thm:integral-apolar-alternatives}による。
この一様な排除は、全ての鉛筆で七次が最初の障害かどうかを必要としない。
\end{proof}

\begin{corollary}[片側が整な階数8の判定]\label{cor:one-integral-apolar-recognition}
隣接階数6・反対階数8の両対称テンソルについて、一方の反対核が整な四次完全交叉を切るとする。
反転四窓が標準ASとなることと、その整曲線上の相異なる次数4の可逆層の対と基点自由鉛筆による
表示を持つことは同値である。このとき両側の反対核は同型な曲線を切り、順序付きデータを復元する。
\end{corollary}
\begin{proof}
標準ASなら定理\ref{thm:all-bisymmetric-as-closure}により滑らかな楕円鉛筆から特殊化する。
反対階数が保たれるので指定側の鉛筆は極限二次式と一致する。整な特異完全交叉四次曲線は
既出の種数公式により節点または尖点を持つ。前定理で非可逆な場合を排除でき、
可逆な場合の十分性・復元は整曲線の鉛筆定理による。
\end{proof}

\begin{proposition}[零次数の成分を許した多重次数判定]\label{prop:semipositive-pencil-balance}
有理曲線の環状多角形、または各枝が射影直線で唯一の特異点が楕円多重点となる曲線について、
$L,M$の多重次数$d,e$を非負とし、各総次数を4とする。
基点自由な$LM$の鉛筆の反転四窓が標準ASとなる必要十分条件は$d=e$である。
\end{proposition}
\begin{proof}
両束とも次数零の成分は切断を取る前に縮められる。環状の鎖の定数を消去すると
両端の値に元の貼り合わせのスカラーの積を課す。一成分だけが残れば二点を結ぶ節点となる。
多重点では零次数枝の定数を共通値とし一次係数を零とすれば、残りの枝に同じ一次の和の条件を得る。
一枝だけなら尖点である。束の積でもこの消去と貼り合わせ条件は可換である。
従って$d=e$なら既出の正次数の場合に帰着する。以下、全成分で$d_i+e_i>0$とする。成分数は高々8である。

相補的な非空の連続部分鎖$I,J$について、$I$だけに支えを持つ$L$の切断の次元は
$\max(d(I)-1,0)$である。鎖の部分の総次数が正なら、零次数の連続部分の定数を消去した後に、
二つの端点の値を独立に選べるためである。総次数零なら端点で消える定数は零である。
楕円多重点では任意の真部分集合で同じ式が成立する。共通値を零とし、
次数が正の枝があれば一次係数に独立な一式を課すからである。
これらの部分に支えを持つ相補的な切断の積は零で、各々の双対への射影で全二次関係が消える。
ゆえに二次元数の積が2以上なら指数成長となる。$d,e$の役割を交換した条件も使える。

この検査に残る$d\ne e$は、成分と束の順序を除き
$((4,0),(3,1))$と$((3,1,0),(3,0,1))$だけである。
有限列挙の方法は、各成分への四単位の配分を二回選び、$d_i+e_i>0$を課し、
全ての巡回区間$J$で和$d(J),e(J)$が同じ区分
$\{0,1\},\{2\},\{3,4\}$に属するものを残すことである。
区分が異なれば上の積の少なくとも一方が2以上となる。
\path{computations/verify_binary_product_obstruction.py}は成分数2から8を厳密列挙する。
多重点では枝に任意の巡回順を付けてよい。

二成分の場合、主成分での第一束の切断のうち、二つの貼り合わせ点
（接する場合はその長さ2の導手）で消えるものは$hP_2$で、他成分で恒等零である。
第二束は主成分へ全三次式として制限できる。従って積は$P_2,P_3$の積を通る。
三成分の場合には両束とも主成分上で全三次式を取れる。環状なら主成分の二端のうち、
第一束が次数零の枝と交わる点において消える三次式を第一束に取る。これは一次式倍の$P_2$である。
第一束の切断はその零次数枝全体でも消える。残りの枝の上では第二束が定数なので、
積は主成分での積の端点の値と、もう一端の値零から一意に決まる。
三重点なら三次式を共通点で消す。第一束は零次数枝で消え、残りの補助枝では第二束が定数である。
そこでの積の値は零で、一次係数は主成分上の積の一次係数から一意に決まる
（一次係数の和の条件による）。いずれも選んだ切断の積は主成分の$P_2\otimes P_3\to P_5$を通る。
補題\ref{lem:binary-diagonal-paths}の商により、残った二型も指数成長する。
\end{proof}
